\section{极大元和极小元}

\begin{definition}{极小元,极大元}{}
	设集合$E$配备偏序$\leq$,$a\in E$.
	\begin{enumerate}
		\item 若$\forall x\in E\left(x\leq a\implies x=a\right)$,则称$a$是$E$(在$\leq$下)的极小元.
		\item 若$\forall x\in E\left(x\geq a\implies x=a\right)$,则称$a$是$E$(在$\leq$下)的极大元.
	\end{enumerate}
\end{definition}

\begin{proposition}{极小元和极大元的对偶性}{}
	设集合$E$配备偏序$\leq$,其对偶为$\leq^{\prime}$,$a\in E$.若$a$是$\leq$下的极小元(相对的,极大元),则$a$是$\leq^{\prime}$下的极大元(相对的,极小元).
\end{proposition}

\begin{example}{}{}
	\begin{enumerate}
		\item 设$\mathcal{P}^{\ast}\left(E\right)$是集合$E$的非空子集组成的集合,配备包含偏序,则该集合的极小元组成的集合是单点集组成的集合.
		\item 设$E,F$是集合,给$\PFunc\left(E,F\right)$配备函数延拓偏序,则极大元组成的集合是$\Hom_{\mathsf{Set}}\left(E,F\right)$.
		\item 给$\Z_{>1}$配备整除偏序,则极小元组成的集合是素数集$\mathbb{P}$.
		\item 给$\R$配备标准偏序,那么既没有极大元,也没有极小元.
	\end{enumerate}
\end{example}







